% !TEX root = ../monograph.tex
\chapter{The ultraviolet limit at fixed volume}
\label{ch12:ultraviolet}

We keep the physical period $R_p$ fixed and let the ultraviolet cutoff
$N$ tend to infinity. The level $m$ has spacing $a_m=D_m^{-1}$;
blocking proceeds from $m+1$ to $m$. The purpose of this chapter is to
compare different cutoffs on one fixed configuration space. Compactness
by itself would give only a subsequence. The additional ingredient is an
estimate for stopping the marked evolution at an intermediate level,
with a constant independent of that level.

The common trajectory of Chapter~\ref{ch10:trajectory} and the
marked estimates of Chapter~\ref{ch11:insertions} supply the stability
and finite-difference estimates used here. We state their precise
forms to distinguish a change of cutoff from comparison of two
insertions under one fixed transformation.

\section{Fixed-level convergence data}
\label{ch12:inputs}

Fix finitely many labels $v_i\in V_{d_i}$ and distinct physical origins.
At a fixed level $m$ above their reference levels, index the activities
by the same local supports and prescribed sets of large-field sites for
every cutoff $N>m$. The small-field indicators and eligibility criteria
are also fixed at level $m$. A large-field profile is the sum of activities
having one particular support and prescribed set of large-field sites; the original activity
list can be infinite even though the number of such profiles on a
fixed finite torus is finite.

The required background convergence consists of Cauchy convergence of
$\beta_{N,p,m}$, of the path coefficients in their exponentially
weighted sums, and of the regular and large-field activity arrays in
the norms of Section~\ref{ch04:activities}. The scalar multiple of the
density need not converge. For the marked transformation we need
fixed-input continuity across cutoffs: at each fixed level, applying
the cutoff-dependent marked map to any fixed admissible remainder
graph gives Cauchy output graphs and coefficient maps. This condition
is stronger than contraction for two graphs acted on by the same map.
For the common construction, Theorem~\ref{ch10:convergence} supplies
the background limits. The marked estimates and fixed-input comparison
in Section~\ref{ch11:one-graph} supply the corresponding single-insertion
maps. We keep the needed forms explicit to show how the two
comparisons enter this argument.

Here is a useful way to obtain the required path convergence. For
canonical path indices $a$ of length $\ell(a)$, suppose the coefficients
$c_{N,a}$ converge individually and obey a common bound at width
$w>v$. Writing $c_a=\lim_{N\to\infty}c_{N,a}$ and $C$ for the common bound, we have
\[
 |c_{N,a}-c_a|e^{v\ell(a)}\le 2C e^{-(w-v)\ell(a)}.
\]
The canonical indices enumerate displacement vectors and finitely many
plaquette directions, rather than all walks of a given length. For a
row their length is $|x|_1+4$, and
\[
 \sum_{x\in\ZZ^4}e^{-\delta(|x|_1+4)}
 =e^{-4\delta}\left(\frac{1+e^{-\delta}}{1-e^{-\delta}}\right)^4<\infty.
\]
For an atom of degree $q$ the analogous sum has $q-1$ displacement
vectors and exponent $4(q-1)$. Dominated convergence therefore gives
convergence in the smaller width $v$. A strict width reserve makes
coordinatewise coefficient convergence sufficient for the actual
path-action norm; convergence of an unrelated finite-volume matrix
would not do so.

\section{Contraction of the complete marked remainder}
\label{ch12:graphs}

For a fixed degree, let $\mathcal R_m$ denote the space of remainder
arrays at level $m$. Its seminorm is the infimum of the sum of the
regular derivative majorants and the large-field value majorants
multiplied by $\kappa_m^{-1}$, where
$\kappa_m=(\epsilon_m^{\rm for})^{1/2}$ is the square root of the
actual level-$m$ large-field bound, as in
Section~\ref{ch11:one-graph}. We keep the arrays themselves, without
identifying different arrays merely because their evaluations agree.
Evaluation of the regular sum and of each large-field profile is a
bounded linear map in this seminorm. This follows first for any
admissible majorant and then by taking its infimum.

An admissible graph $G:V_d\to\mathcal R_m$ obeys
$\|Gx\|_{\mathcal R_m}\le s_m\|x\|_1$, where $s_m\ge1$ is fixed
by the finite-degree construction. Set
\[
 d_m(G,H)=\sup_{x\ne0}
       \frac{\|(G-H)x\|_{\mathcal R_m}}{s_m\|x\|_1}.
\]
Thus any two admissible graphs have distance at most two. A marked step
from $m+1$ to $m$ sends the input $(x,Gx)$ to $(I_Gx,J_Gx)$, where
$I_Gx$ is the new polynomial coefficient and $J_Gx$ the new remainder.
Using the new coefficient as independent variable gives the output graph
$\Phi G=J_GI_G^{-1}$.

\begin{lemma}[Graph contraction]
\label{ch12:contraction}
Suppose the marked step maps admissible graphs into admissible graphs,
and the common-support marked difference estimates give, for
$a=d_{m+1}(G,H)$,
\[
 \|I_G-I_H\|\le\delta_m a,\qquad
 \|(J_G-J_H)x\|_{\mathcal R_m}\le\delta_m a\|x\|_1,
 \qquad \|I_G^{-1}\|\le s_{m+1}.
\]
If $4\delta_m s_{m+1}\le1$, then
$d_m(\Phi G,\Phi H)\le a/2$.
\end{lemma}
\begin{proof}
For $y=I_G^{-1}x$, use $J_H=(\Phi H)I_H$ to write
\[
 (\Phi G-\Phi H)x=(J_G-J_H)y+(\Phi H)(I_H-I_G)y.
\]
Its norm is at most $(1+s_m)\delta_m a s_{m+1}\|x\|_1$,
which is at most $as_m\|x\|_1/2$.
\end{proof}
The difference estimates are for the complete regular and large-field
remainder on common supports. Linearity of a single insertion
allows its local polynomial part to cancel. A contraction only for the regular
part would not prove this lemma for the actual inserted density.
For the actual transformation, \eqref{ch11:graph-estimates} gives
ball invariance and the inverse bound, while
\eqref{ch11:graph-difference} gives these two differences with
$\delta_m=C_d\epsilon_m s_{m+1}$.
The selected ratio and thresholds for the starting scales impose
$4\delta_m s_{m+1}\le1$. The fine and coarse remainder norms retain
their distinct factors $\kappa_{m+1}^{-1}$ and $\kappa_m^{-1}$.
Weighted path continuity at a strict width reserve and activity-norm
continuity in Section~\ref{ch11:one-graph}, together with
Theorem~\ref{ch10:convergence}, give the fixed-input cutoff
continuity as well.

\begin{proposition}[Independence of the terminal graph in the cutoff limit]
\label{ch12:forgetting}
For the selected trajectory, the graph contraction and fixed-input
cutoff continuity just verified imply that the selected
graphs and coefficients normalized at one reference level are Cauchy
at every fixed level.
\end{proposition}
\begin{proof}
Fix $m$ and a positive integer $h$. Replace the selected graph at
$m+h$ by one fixed admissible graph and propagate down $h$ steps.
The resulting graphs are Cauchy as $N\to\infty$: fixed-input
continuity and contraction give this by finite induction. Their
distance from the actual graphs is at most $2^{1-h}$. First let $h$
be large and then let both cutoffs tend to infinity. This proves the
Cauchy assertion without an infinite product of inverse maps.

The coefficient at its reference level is the prescribed label.
For each subsequent fixed level solve the finite-dimensional relation
$x_{m,N}=I_{m,N}x_{m+1,N}$. The maps are Cauchy and their inverses
uniformly bounded at this fixed level. The identity
\[
 \|I_{m,N}^{-1}x_{m,N}-I_{m,M}^{-1}x_{m,M}\|_1
 \le C_m\bigl(\|x_{m,N}-x_{m,M}\|_1
       +\|I_{m,N}-I_{m,M}\|\,\|I_{m,M}^{-1}x_{m,M}\|_1\bigr)
\]
proves the assertion by finite induction. Evaluation of $G_{m,N}x_{m,N}$
is Cauchy by the graph bound and its bounded evaluation maps.
\end{proof}

Consequently the complete functions multiplying a single source $f_{i,N,m}(U)$ and
large-field profiles linear in the sources $k_{i,N,m,Q}(U)$ converge uniformly in
$U$. Here $Q$ is a support together with its exact prescribed set of large-field sites.
Multiplication by a fixed Borel event indicator preserves uniform convergence;
no continuity of the event indicator as a function of $U$ is required.

\section{The scalar normalization cancels}
\label{ch12:normalization}

Write $\rho_{N,m}$ for the actual positive density at the fixed period,
$s_{N,m}$ for its extensive scalar term, and
\[
 \widetilde\rho_{N,m}=e^{-s_{N,m}}\rho_{N,m},\qquad
 Z^0_{N,m}=\int\widetilde\rho_{N,m}\,d\mu_{m,p},
\]
where $\mu_{m,p}$ is product probability Haar measure. Background
convergence makes $\widetilde\rho_{N,m}$ uniformly Cauchy. Indeed,
curvatures are bounded on the compact group, so weighted path
convergence implies uniform action convergence. Summing a rooted
activity norm over the finitely many sites gives uniform convergence
of the regular sum and every large-field profile. The exact-cover gas
is a finite polynomial in the latter profiles and fixed eligibility
indicators. Its product with the exponential of the regular action
therefore converges uniformly.

The limit of $Z^0_{N,m}$ is strictly positive. Choose a fixed link ball
$\mathcal U_m$ about the identity small enough that every inspected
plaquette satisfies the level-$m$ cutoff. On this set there are no
large-field sites, so the exact-cover sum is one. The action and the
regular sum are bounded above by a constant $A_{m,p}$ uniformly on
the cutoff tail. Since Haar measure assigns positive mass to every
nonempty open link ball,
\begin{equation}
 Z^0_{N,m}\ge e^{-A_{m,p}}\mu_{m,p}(\mathcal U_m)>0.
 \label{ch12:positive-partition}
\end{equation}
Nonnegativity of the full density justifies restricting the integral
to $\mathcal U_m$. Positivity of each partition function alone would
not imply a positive limit.

Let $z_1,\ldots,z_n$ be commuting variables with $z_i^2=0$, and write
$z_S=\prod_{i\in S}z_i$. The exact marked density at level $m$ has
regular terms linear in the sources $\sum_i z_i f_i$, higher regular entries
$H_X$, and large-field factors
\[
 K_Q+\sum_i z_i k_{i,Q}+K_Q^{\ge2}.
\]
Every term in $H_X$ and $K_Q^{\ge2}$ contains at least two distinct
source variables. All insertion origins belong to its support. Keep the complete
terms linear in the sources and set only these higher entries to zero. Call the
resulting density $\mathscr X^0_{N,m}(z,U)$. The expression
\begin{equation}
 C_{N,p,m}(v;x)=
 \frac{\int[z_{\{1,\ldots,n\}}]
            e^{-s_{N,m}}\mathscr X^0_{N,m}(z,U)\,d\mu_{m,p}}
      {Z^0_{N,m}}
 \label{ch12:stopped}
\end{equation}
is Cauchy in $N$ for fixed $m$. Its numerator is a finite polynomial
in uniformly convergent grouped profiles and functions multiplying a single source,
multiplied by a uniformly convergent exponential. Its denominator has
the positive limit in \eqref{ch12:positive-partition}. The same scalar
has been removed from both. In particular, deleting terms involving at least two insertions does
not mean deleting the large-field profiles linear in the sources.

\section{A stopping estimate uniform in the stopping level}
\label{ch12:stopping}

Recall that $B_m$ is the support radius from \eqref{ch04:thresholds}
and $\operatorname{cov}_{B_m}(X)$ counts a connected cover of $X$ by
radius-$B_m$ boxes. Suppose the physical origins belong to one common
lattice and are distinct modulo $R_p$. Their separation at level $m$
is at least $\eta/a_m$ for some $\eta>0$. Since $a_mB_m^2\to0$,
choose a fixed base level $r$ so that separation is at least $B_m^2$
for every $m\ge r$. Enlarge $r$ to include all reference levels.

We use the following marked norm conclusion. Terms involving at least two insertions have rooted
majorant sums at most one at support exponent $\sigma_m$. The transfer
accepts exponent $\sigma_{m+1}-4$ and restores exponent $\sigma_m$.
It retains all inventories and insertion origins. Moreover, after transfer
to the fixed base level, normalized coefficient extraction is bounded
by
\begin{equation}
 C_r=\frac{\exp\bigl(|\Lambda_{r,p}|q_r^{\rm site}+n(x_r+y_r+2)\bigr)}
                 {\prod_i w_i}.
 \label{ch12:root-bound}
\end{equation}
Here $q_r^{\rm site}$ is the per-site normalization cost in \eqref{ch11:replicated-bound}; $w_i>0$ are fixed source weights; $x_r$ bounds the weighted regular sizes of the terms linear in the sources; and $y_r$
bounds their weighted large-field profile sums. Such weights can be
chosen uniformly on the cutoff tail: the fixed-level coefficients linear in the sources are bounded, so scale each weight down to its positive
allowed size. Crucially all four quantities in \eqref{ch12:root-bound}
are evaluated at $r$, not at the later stopping level $m$.
These statements concern the complete actual marked transfer.
They follow from the estimate for terms involving at least two insertions
\eqref{ch11:higher-induction} and the normalized chessboard bound
\eqref{ch11:root-bound}.

\begin{lemma}[Removal of terms involving at least two insertions]
\label{ch12:omission}
Under these marked norm and transfer bounds, for each $m>r$ and all
sufficiently large $N$,
\[
 |M_{N,p}(v;x/a_N)-C_{N,p,m}(v;x)|
       \le 2C_r e^{-B_m/64}.
\]
The constant $C_r$ is independent of $m,N$; the cutoff threshold can
depend on $m$.
\end{lemma}
\begin{proof}
A support containing two origins separated by $B_m^2$ has cover
number at least $B_m/64$. Multiply every term of source degree at least two, and only those
entries, by a complex scalar $u$. For $|u|<S_m=e^{B_m/64}$ its
majorant still fits exponent $\sigma_m-2$, since
\[
 |u|e^{(\sigma_m-2)\operatorname{cov}_{B_m}(X)}
       \le e^{\sigma_m\operatorname{cov}_{B_m}(X)}.
\]
Transport the entire modified marked density to $r$. The support
reserve allows its first and subsequent steps, and the constant and
coefficients linear in the sources are unchanged. If a source monomial contains $q$ origins,
the stronger geometric bound $B_mq\le64\operatorname{cov}_{B_m}(X)$
also absorbs a change of each source weight by at most $e^{B_m/128}$.
This is the weight reserve needed when using the marked transfer.

Let $F_{N,m}(u)$ be the normalized integral of its full-source
coefficient. Exact probability-kernel integration preserves that
integral and its vacuum denominator. The fixed-base estimate gives
$|F_{N,m}(u)|\le C_r$ throughout the disk. The source algebra is
nilpotent; after grouping the finitely many support and prescribed-site profiles,
$F_{N,m}$ is a polynomial in $u$, hence holomorphic. For $S_m>2$,
Cauchy's derivative estimate on circles of radius $S_m/2$ about points
of $[0,1]$ gives
$|F_{N,m}(1)-F_{N,m}(0)|\le2C_r/S_m$. These two values are the
actual cutoff moment and \eqref{ch12:stopped}, respectively.
\end{proof}

\begin{proposition}[Cauchy convergence on common lattice points]
\label{ch12:dense-cauchy}
For the selected trajectory, the correlations converge along the whole cutoff sequence at every
configuration of distinct points belonging to one common lattice.
\end{proposition}
\begin{proof}
Given an error tolerance, first choose $m$ so that the two stopping
errors in Lemma~\ref{ch12:omission} are small. Keep $m$ fixed. The
stopped expression \eqref{ch12:stopped} is Cauchy as both cutoffs tend
to infinity. The triangle inequality proves the assertion. No new
subsequence is chosen.
\end{proof}

\section{Smoothness and all lattice differences}
\label{ch12:smooth}

The remaining passage is a standard discrete version of the
Arzel\`a--Ascoli theorem and convergence of derivatives; compare
\cite[Theorems~7.17 and~7.25]{Rudin1976}. We record the precise
application, since the rounded lattice functions are not continuous.
Let $U\subset\RR^D$ be open, $a_N\downarrow0$, and let $F_N$ be
functions on $a_N\ZZ^D$. Write
\[
 \Delta_{N,j}F_N(x)=a_N^{-1}(F_N(x+a_Ne_j)-F_N(x)).
\]
At a point $x\in U$, a rounded value means evaluation at a lattice point
within one mesh step in each coordinate, chosen by a fixed rule. Assume every fixed iterated difference is bounded on each compact
subset of $U$, eventually in $N$, and the rounded values are Cauchy
on a dense subset. Then there is $F\in C^\infty(U)$ such that all
rounded iterated differences converge locally uniformly to the
corresponding derivatives of $F$.

Only the discrete-to-continuous adaptation needs explanation. For
nearby lattice points, telescoping along coordinate edges gives
$|G_N(y)-G_N(x)|\le C|y-x|_1$ whenever the first differences of $G_N$
are bounded on a slightly larger compact neighborhood. Rounding adds
at most $2DCa_N$ to this modulus. A finite net from the dense Cauchy
set therefore proves locally uniform convergence of $F_N$.
Bounded second differences give the discrete Taylor estimate
\[
 |G_N(y)-G_N(x)-\sum_j(y_j-x_j)\Delta_{N,j}G_N(x)|
       \le C|y-x|_1^2.
\]
For fixed small $h>0$, compare $\Delta_{N,j}G_N(x_N)$ to the secant
with interval $h_N=a_N\lceil h/a_N\rceil$. The error is at most
$Ch_N$. The secants converge first as $N\to\infty$; then $h\to0$
proves that the first differences are Cauchy. Their own first-difference
bounds give locally uniform convergence. Passing to the displayed
Taylor inequality identifies their limits as the derivatives of the
limit of $G_N$. Iteration proves the assertion for every finite order.
All paths and shifted stencils lie in the enlarged compact neighborhood.

Apply this to periodic lifts on
\[
 U_{p,n}=\{(x_1,\ldots,x_n)\in(\RR^4)^n:
                         x_i-x_j\notin R_p\ZZ^4\text{ for }i\ne j\}.
\]
The union of the nested physical lattices is dense here. Its points
supply Proposition~\ref{ch12:dense-cauchy}; the all-order separated
finite-difference bounds supply the other hypothesis. Points sharing
one or several coordinate projections are allowed. Only coincidences
of the full four-dimensional points modulo the period are excluded.
The limiting function, denoted $M^{\rm UV}_{p,v}$, is smooth and
periodic on this domain.
The difference hypothesis is supplied by
Theorem~\ref{ch11:all-period-differences}: on a compact subset
of $U_{p,n}$ the periodic separation has a positive lower bound,
and any fixed shifted difference stencil stays in a slightly
larger separated compact set on the cutoff tail.

\section{The fixed-volume conclusion}
\label{ch12:conclusion}

\begin{theorem}[Full ultraviolet limit]
\label{ch12:uv-theorem}
For the common selected trajectory with the strengthened stage
requirements of Chapters~\ref{ch10:trajectory} and
\ref{ch11:insertions}, the correlations have the smooth
full-sequence limit
$M^{\rm UV}_{p,v}$ at every fixed physical period. For every
$f\in C_c^\infty(U_{p,n})$ supported in an admissible period cell
and every finite coordinate list $\alpha$,
\[
 \lim_{N\to\infty}a_N^{4n}\sum_{y}
    f(a_Ny)\Delta_{N,\alpha}M_{N,p}(v;y)
       =\int f(x)\partial^\alpha M^{\rm UV}_{p,v}(x)\,dx.
\]
Here differences in the summand are divided by their physical step
lengths. Their shifted sites must lie in the slightly enlarged
separated set, which holds on a cutoff tail.
\end{theorem}
\begin{proof}
The preceding section gives the full smooth limit on separated configurations.
The last assertion is a Riemann-sum consequence of locally uniform
convergence. Fixed forward, backward or centered stencils give the
same limit: telescope each stencil as a finite average of shifted
unit forward differences. The averaging weights sum to one and the
shifted points approach the same point.
\end{proof}

The convergence threshold in this theorem may depend on $p$. Bounds
on compact separated sets that were uniform in $p$ remain uniform
for the limiting kernels, by passage to the limit. These two
quantifier statements serve different purposes: the former removes
the ultraviolet cutoff; the latter permits the subsequent
infinite-volume extraction. Finite reflection inequalities also pass
to the limit whenever every correlation in their finite quadratic
form is supported away from collisions. Neither this argument nor
reflection positivity asserts an unsubtracted limit at coincident
points. Distributional extensions and reconstruction require their
own arguments.
